\section{Conclusions}
\label{vii:sec:conclusiones}

L'amplitude torsionnelle est déterminée par l'action spinorielle et
l'état matériel, avec une formule explicite et un secteur dans lequel sa
positivité est conservée. Ce résultat permet d'obtenir une contribution
énergétique proportionnelle à \(-a^{-6}\), un rebond homogène unique sous les
conditions démontrées et, pour une poussière plate avec
\(\Lambda_{\mathrm{eff}}=0\), une évolution régulière pour tout temps avec
des géodésiques causales de Levi--Civita complètes. La séquence relie un
observable quartique de l'état à une échelle cosmologique minimale. Son
origine demeure dans la construction APP--TRIT--TPK : l'orientation, le résidu,
la retenue et la profondeur sont transportés jusqu'aux actions, aux courants et aux
observables de la réalisation géométrique.

\subsection{Détermination de l'amplitude et évolution cosmologique}

Le résultat central de la réalisation spinorielle est la détermination
de l'amplitude torsionnelle comme fonctionnelle de l'état matériel :
\[
 s_0[\varrho,V_c]
 =\frac{3\kappa c^2\hbar^2}{16}
   \frac{\gamma^2}{1+\gamma^2}
   \frac{\operatorname{Tr}
    (\varrho\widehat{\mathscr Q}_{\rm sp})}{V_c^2}.
\]
La formule provient successivement de la représentation interne de Clifford, du courant de spin, de l'élimination de la contorsion, de l'observable quartique normalement ordonné et des variations du lapse et du volume. Le facteur quadratique du volume exprime la dépendance de la densité à l'égard du second moment du courant. Une moyenne de spin nulle conserve un second moment qui peut contribuer à la géométrie.

Dans le secteur d'occupation trois, on obtient
\(\widehat{\mathscr Q}_{\rm sp}=4I\). Pour des états de trace un
à support dans ce secteur, toute évolution unitaire qui le préserve
conserve son espérance et détermine une amplitude positive.
L'état matériel et la cellule comobile individualisent la réalisation ;
\(G\), \(\hbar\) et \(\gamma\) conservent leurs antécédents
structurels. La constante de Planck réduite intervient quadratiquement dans
l'interaction, Barbero--Immirzi détermine le coefficient de l'
opérateur de Holst et la gravitation transforme cette interaction en
réponse géométrique. Cette composition explique conjointement
la valeur fonctionnelle de l'amplitude et la fonction de chaque facteur.

La loi de mémoire \(M_n=\varphi^{-2n}\) et la lecture spatiale
\(a_n/a_{\rm ref}=\varphi^{n/3}\) produisent la dilution
\(a^{-6}\). Dans la réalisation spinorielle à second moment
conservé, cette même puissance apparaît par la variation matérielle.
Avec \(\rho_0,s_0>0\), \(w<1\), une courbure spatiale nulle et
\(\Lambda_{\rm eff}=0\), le seuil de rebond est unique et
transverse. Deux bornes différenciées préservent localement son
existence, son unicité et son orientation temporelle sous perturbations.

Le secteur de poussière admet en outre une solution exacte pour tout temps. Le facteur d'échelle reste séparé de zéro, la courbure est finie et les géodésiques causales de Levi--Civita sont complètes dans les deux directions. Le résultat global appartient à cette réalisation explicite ; la persistance du rebond est établie sous les conditions locales de son théorème. Les deux preuves montrent comment l'amplitude obtenue à partir de l'état se traduit en une échelle minimale et en une évolution de contraction et d'expansion.

\subsection{De la mémoire transportée à la réponse gravitationnelle}

La normalisation radiale s'obtient après la composition longitudinale
\(L_\alpha=(5759/23040)\alpha^{16}L_*\). Le transport normalisé
de l'archive et la métrique positive produisent
\(R_X=L_\alpha\mathsf U X\mathsf U^*\) ; l'énergie conserve le même
caractère et sa longueur conjuguée satisfait
\(H_X\Lambda_X=\hbar_{\rm ret}cI\). Dans la lecture de rayon réduit,
ces égalités déterminent \(G=c^3L_\alpha^2/\hbar_{\rm ret}\), commun
aux modes positifs de la réalisation. La preuve conserve ses règles
longitudinale et métrique et l'identification physique déclarée. La réduction
conjointe de la frontière et de la mémoire maintient le coefficient ; dans la réduction
dynamique, on transporte aussi le paramètre spectral et la norme temporelle.
L'expression circulaire est retrouvée avec
\(t_0=\pi_{\rm HMT}L_\alpha/(54c)\).

L'espace à cinq composantes s'obtient par une projection explicite
des douze sommets et est transporté isométriquement vers des tenseurs
symétriques sans trace. La réduction transverse conserve deux composantes.
La séquence \(12\to5\to5\to2\) distingue l'incidence initiale,
le porteur gravitationnel, sa représentation tensorielle et la lecture
transverse. Le coefficient \(G_a\) fixe le couplage commun ;
les orientations appartiennent à l'opérateur et à ses composantes.

Dans la même section de retour, le coefficient de l'amplitude spinorielle admet l'expression équivalente
\[
 s_0[\varrho,V_c]
 =\frac{3\pi c\hbar_{\rm ret}L_\alpha^2}{2}
   \frac{\gamma^2}{1+\gamma^2}
   \frac{\operatorname{Tr}(\varrho\widehat{\mathscr Q}_{\rm sp})}{V_c^2}.
\]
Elle résulte de la substitution de \(\kappa=8\pi L_\alpha^2/(\hbar_{\rm ret}c)\) dans la fonctionnelle déjà démontrée. La longueur et l'action fixent ainsi son coefficient, tandis que l'état et le volume conservent leurs fonctions matérielles dans la solution cosmologique.

L'écran de mémoire fournit une soudure et une réponse
énergétique compatibles avec le raffinement. Son extension minimisante,
construite à partir de l'incidence transportée, détermine une
énergie et une différentielle par rapport à la connexion. Le calcul inclut
les variations de l'incidence, de la frontière, du poids et de la normalisation. Lors du
passage à la réalisation de Cartan, la cotétrade et l'appariement
matériel convertissent cette différentielle en un courant de domaine
géométrique explicite.

Les équations de connexion admettent une inversion algébrique des opérateurs de Holst et de torsion. Dans le secteur matériel affine par rapport à la contorsion, cette inversion produit la correction effective de l'action. Pour l'extension minimale, dont le courant dépend en général de la connexion elle-même, on démontre l'équation stationnaire complète, le critère local de continuation au moyen de sa hessienne et la formule intégrale de l'action effective. L'action spinorielle et l'action d'extension minimisante conservent donc leurs arguments et leurs lois de variation propres. La première fournit la fonctionnelle d'amplitude évaluée dans la solution cosmologique ; la seconde décrit la réponse variationnelle complète de l'écran transporté.

\subsection{Information, horizons et réponse relationnelle}

Les bilans résiduels réunissent dans une même formulation la trace, la partie sans trace et leurs divergences. Avec des couplages constants dans la coordinatisation et une source matérielle conservée, le tenseur résiduel conserve covariamment son bilan total ; la variation de sa trace détermine l'échange entre les deux secteurs. La constance de la trace caractérise la conservation séparée.

L'information d'horizon utilise l'unité surfacique,
l'action et la conversion thermique préalablement construites.
L'état réduit et sa purification distinguent l'information
accessible et les corrélations conservées. Dans le bilan sphérique
développé se vérifient conjointement \(TS=E\),
\(T\,dS=2\,dE\) et \(S\,dT=-dE\) ; la dépendance de
la température vis-à-vis du rayon fait partie de ces identités.
La constante de Boltzmann convertit le contenu informationnel
dans la section thermique, tandis que la relation action--gravitation
détermine l'unité d'aire de sa lecture géométrique.

\Needspace{24\baselineskip}
La composition aréale et chronologique détermine
\(\mathcal T_{I/A}q_{\mathcal A}/\log3=h/(108t_0\log3)\).
En \(R=L_\alpha\), les températures cosmologique et chronologique
sont reliées par le facteur \(\log3/(2\pi)\).
Pour le hamiltonien de réponse
\(\mathsf H_R=\tau_R\Delta_0\mathsf D_\partial\), les poids
90 et 120 correspondent à des sauts énergétiques distincts d'échelle
commune. Le bilan fini est exactement
\(\mathcal E_R-T_{\rm cos}\mathcal S_R
=\tau_R\mathcal D(\sigma\Vert\rho_0)\), et son terme tangent
est la première loi à rayon fixé. La section relative construite à partir de
cet état détermine un unique transducteur positif ; l'identification
d'un lecteur supplémentaire se réduit à vérifier sa valeur sur la section
chronologique de~\eqref{viii:eq:criterio-una-seccion}. Les isométries qui
conservent le registre transportent ce bilan et ses observables.

L'observateur est défini au moyen d'opérations d'orientation, de trivialisation et de restriction des observables. La réponse auto-inertielle conserve sa factorisation par l'information globale et explicite le défaut produit par la projection. Enfin, dans le canal chronologique réalisé, la différence covariante construite avec le transport des routes admet, pour des sources à support fini, des opérateurs de Green retardés et avancés. La compatibilité des deux lectures à la frontière délimite un espace de sources et une projection qui intervient dans la réponse variationnelle. La bidirectionnalité est exprimée par des opérateurs, des moments et des conditions de frontière qui peuvent être vérifiés dans l'article.

\enlargethispage{2\baselineskip}
La relation décisive est que l'information conservée par l'état
intervient dans la réponse géométrique au moyen d'un courant et de son second
moment. L'élimination de la contorsion et la variation métrique convertissent
cet observable en une densité et une pression ; la solution de poussière montre
comment toutes deux déterminent une échelle minimale et une évolution causale complète.
Les bilans et la lecture de l'observateur conservent à leur tour la
distinction entre information accessible et corrélation globale. Dans les
domaines construits, la mémoire, la source matérielle et la géométrie forment ainsi une
chaîne explicite de déterminations, du noyau arithmétique jusqu'à
l'évolution gravitationnelle et son contenu informationnel.
